\cotefile{33}
\scanpage{147}
\margin{\unsure{146}}
\mspage{10.}

Alors le topos minimal cherché sera $\operatorname{Isloc}(X)$ et le théorème sera démontré. Pour cela, associons à chaque $F\inserted{\in}\overset{\inserted{\sim}}{\operatorname{Ouv}}(X)$ l'espace topologique étalé au-dessus de $X$ dans le sens classique (cf.\ \inserted{[}6\inserted{]},II.1.2 ) et munissons cet espace de la structure annelée induite par celle de $X$. Nous obtenons ainsi un objet de $\operatorname{Isloc}(X)$ et il est facile de voir que le foncteur ainsi construit est une équivalence de catégories. Le théorème 1 est donc démontré.

\uline{Remarque 1.}- La condition essentielle dans le théorème 1 c'est (A)
La condition (B) est seulement une condition de régularité qui assurera que les $X$-préschémas en questions soient \uline{localement de présentation finie sur} $X$ (cf.\ \inserted{[}EGA IV.8.14\inserted{]} et \inserted{[}SGA 1963 Exp.\ $\quad$\inserted{]} ).

\uline{Remarque 2.}- En vertu des résultats du \S\ 1, on peut interpréter le théorème 1 pour (Sch) en disant que la topologie de Zariski sur (Sch) est la topologie la plus fine qui ``ignore les faisceaux structuraux des préschémas et ne s'intéresse qu'à leurs espaces topologiques sous-jacents''. Mais il est évident qu'on ne pourra pas se contenter d'étudier une topologie aussi faible, ne serait-ce que parce que tout préschéma $\operatorname{Spec}(k)$, où $k$ est un corps, aurait la cohomologie nulle en dimension $\inserted{>}0$, et l'on ne pourrait donc pas espérer de distinguer les corps entre eux par des moyens cohomologiques si l'on se restreignait à la topologie de Zariski. D'autre part, on s'est aperçu depuis longtemps qu'il existe des fibrés algébriques intéressants qui ne sont pas localement triviaux pour la topologie de Zariski, mais qui devraient l'\inserted{ê}tre dans un certain sens\inserted{(voir [}4\inserted{]} et \inserted{[}7\inserted{])}, C'est d'ailleurs par inspiration de ces travaux de Serre \inserted{[}4\inserted{]} sur des fibrés algébriques que Grothendieck a été amené à introduire et à étudier des topologies plus fortes que (Zar) sur (Sch)
